% !TEX root = ../../../main.tex
% Block card: detection.symbol_search -- "Symbol search" (Part I, the blocks chapter).
\cardfield{What it does}
Searches the letter string made by a symbolic encoding for a pattern you write, and marks
the stretch of time behind each match.

\cardfield{When to use it}
After \blockref{detection.sax_dsax} (or another block that turns a span into letters), when
you can describe the shape you want in letters: ``one or more rises then one or more falls''
is \texttt{c+a+} at dSAX's default three letters (\emph{a} down, \emph{b} same, \emph{c}
up), and that is the default pattern.

\begin{plainterms}[Patterns (regular expressions)]
A pattern is a small recipe for matching text, like a search box that understands
``one or more'' and ``any of''. \texttt{c} matches the letter \emph{c};
\texttt{c+} matches one or more \emph{c}s in a row; \texttt{a\{3,\}} matches three or more
\emph{a}s; \texttt{[ab]} matches either letter; \texttt{b*} matches any number of \emph{b}s,
none included. So \texttt{c+b*a\{3,\}} reads ``a rise, perhaps a pause, then a fall lasting
at least three pieces''.
\end{plainterms}

\cardfield{What goes in, what comes out}
In goes the encoding from the step before. Out come spans: one per match, starting where
the first matched letter's piece of time starts and ending where the last one's ends. Each
span is labelled with the letters it matched.

\cardfield{Settings that matter}
\emph{pattern} is the search. \emph{alphabet} changes the letters used to spell the
symbols, for encodings with their own (\texttt{dDSUu} for the five-stage encoding); leave it
empty for the SAX blocks. \emph{overlapping} lets a match start inside the previous one.
The block page tells you before a run whether the pattern is valid; it does not count
matches until you run.

\cardfield{Watch out for}
A letter means different things after different encoders: \emph{c} is ``rising'' after dSAX
but ``the third level'' after cSAX or pSAX. A pattern that can match nothing at all, such
as \texttt{c*}, can stop the search early and return no spans. Encodings with more than 26
letters spell some symbols with two characters, and the spans then land in the wrong place.
